% ============================================================
% DOCUMENTATION COMPLÈTE
% Système de Gestion des Ressources Matérielles (S-GRM)
% ------------------------------------------------------------
% Livrable 1 : Cahier des Charges Fonctionnel (CDC)
% Livrable 2 : Document d'Architecture Technique (DAT)
% ============================================================
% Compilation : pdflatex (au moins 2 passes pour la TOC)
%               + plantuml externe pour générer les images
%               (les codes PlantUML sont fournis en lstlisting)
% ============================================================

\documentclass[a4paper,11pt,openany]{report}

% --- Encodage et langue ---
\usepackage[utf8]{inputenc}
\usepackage[T1]{fontenc}
\usepackage[french]{babel}
\usepackage{lmodern}
\usepackage{microtype}

% --- Mise en page ---
\usepackage[a4paper,top=2.5cm,bottom=2.5cm,left=2.7cm,right=2.7cm]{geometry}
\usepackage{setspace}
\onehalfspacing
\usepackage{parskip}

% --- Tableaux et listes ---
\usepackage{array}
\usepackage{longtable}
\usepackage{booktabs}
\usepackage{tabularx}
\usepackage{multirow}
\usepackage{enumitem}

% --- Couleurs et boîtes ---
\usepackage{xcolor}
\definecolor{bleuFac}{RGB}{26,54,93}
\definecolor{bleuClair}{RGB}{43,108,176}
\definecolor{grisCode}{RGB}{245,247,250}
\definecolor{grisFonce}{RGB}{74,85,104}

% --- Code source / PlantUML ---
\usepackage{listings}
\lstset{
  basicstyle=\ttfamily\footnotesize,
  backgroundcolor=\color{grisCode},
  frame=single,
  rulecolor=\color{grisFonce},
  breaklines=true,
  breakatwhitespace=true,
  showstringspaces=false,
  numbers=left,
  numberstyle=\tiny\color{grisFonce},
  stepnumber=1,
  tabsize=2,
  captionpos=b,
  keywordstyle=\color{bleuClair}\bfseries,
  commentstyle=\color{grisFonce}\itshape,
  stringstyle=\color{bleuFac},
  morekeywords={start,stop,if,then,else,endif,actor,participant,
    boundary,control,entity,database,node,component,package,class,
    abstract,extends,implements,interface,note,activate,deactivate,
    return,end,split,fork,merge,detach,partition,group,alt,opt,
    loop,parallel,break,destroy,create,ref,over,right,left,of,
    use,extend,include,as,@startuml,@enduml,@startsalt,@startmindmap,
    skinparam,title,header,footer,legend,top,bottom,center}
}

% --- Liens hypertextes ---
\usepackage[colorlinks=true,linkcolor=bleuFac,
            urlcolor=bleuClair,citecolor=bleuClair]{hyperref}

% --- Mise en forme des titres ---
\usepackage{titlesec}
\titleformat{\chapter}[display]
  {\normalfont\Huge\bfseries\color{bleuFac}}
  {\chaptertitlename\ \thechapter}{20pt}{\Huge}
\titleformat{\section}
  {\normalfont\Large\bfseries\color{bleuFac}}{\thesection}{1em}{}
\titleformat{\subsection}
  {\normalfont\large\bfseries\color{bleuClair}}{\thesubsection}{1em}{}

% --- En-têtes et pieds de page ---
\usepackage{fancyhdr}
\pagestyle{fancy}
\fancyhf{}
\fancyhead[L]{\small\itshape S-GRM — Faculté}
\fancyhead[R]{\small\itshape\nouppercase\leftmark}
\fancyfoot[C]{\thepage}
\renewcommand{\headrulewidth}{0.4pt}

% --- Macros utilitaires ---
\newcommand{\ressource}{\textsc{Ressource}}
\newcommand{\fournisseur}{\textsc{Fournisseur}}
\newcommand{\constat}{\textsc{Constat de Panne}}
\newcommand{\mad}{MAD}

\begin{document}

% ============================================================
% PAGE DE GARDE
% ============================================================
\begin{titlepage}
  \centering
  \vspace*{1.5cm}
  {\Large\bfseries Faculté des Sciences --- Université Privée \par}
  \vspace{0.4cm}
  {\large Filière Master 1 --- Architecture Logicielle\par}
  \vspace{2cm}

  \rule{\linewidth}{1pt}
  \vspace{0.4cm}
  {\Huge\bfseries\color{bleuFac}
    Système de Gestion \\[0.3cm] des Ressources Matérielles\par}
  \vspace{0.4cm}
  \rule{\linewidth}{1pt}

  \vspace{1.5cm}
  {\LARGE\itshape Documentation Technique et Conceptuelle\par}
  \vspace{0.5cm}
  {\Large Cahier des Charges Fonctionnel \\
          \& \\
          Document d'Architecture Technique\par}

  \vfill
  \begin{tabular}{rl}
    \textbf{Auteur}        & : Équipe projet S-GRM \\
    \textbf{Encadrant}     & : Professeur d'Architecture Logicielle \\
    \textbf{Année}         & : 2025--2026 \\
    \textbf{Version}       & : 1.0 \\
    \textbf{Statut}        & : Document validé \\
  \end{tabular}

  \vspace{1cm}
  {\large\today\par}
\end{titlepage}

% ============================================================
% TABLE DES MATIÈRES
% ============================================================
\tableofcontents
\listoftables
\listoffigures

% ============================================================
% RÉSUMÉ EXÉCUTIF
% ============================================================
\chapter*{Résumé exécutif}
\addcontentsline{toc}{chapter}{Résumé exécutif}

Le présent document constitue la documentation officielle du
\textbf{Système de Gestion des Ressources Matérielles} (S-GRM)
destiné à la \textbf{Faculté}. Le système couvre l'intégralité du
cycle de vie du matériel informatique : depuis l'expression
des besoins par les chefs de département, jusqu'à la mise au rebut
des équipements défaillants, en passant par la procédure d'appel
d'offres avec les fournisseurs, l'attribution physique des
matériels (ordinateurs et imprimantes) aux entités utilisatrices,
et enfin le suivi rigoureux des opérations de maintenance.

Cette documentation est structurée en deux parties principales.
La première partie correspond au \textbf{cahier des charges
fonctionnel} ; elle établit le contrat fonctionnel entre la
maîtrise d'ouvrage (la Faculté) et la maîtrise d'œuvre (l'équipe
projet) en formalisant les exigences métier, les acteurs, les
processus et les cas d'utilisation. La seconde partie correspond
au \textbf{Document d'Architecture Technique} (DAT) ; elle
présente les choix techniques retenus pour la mise en œuvre du
système, justifie l'adoption d'une architecture trois-tiers
fondée sur la stack \textsc{PERN} (PostgreSQL, Express.js, React,
Node.js), et détaille la modélisation orientée-objet sous forme de
diagrammes UML.

L'ensemble du présent document a été rédigé selon les standards
académiques en vigueur. Les diagrammes sont fournis en notation
\textbf{PlantUML}, ce qui permet leur génération automatisée et
leur révision collaborative dans un système de gestion de versions.

\part{Cahier des Charges Fonctionnel}

% ============================================================
% CHAPITRE 1 — CONTEXTE & ACTEURS
% ============================================================
\chapter{Contexte du projet et description des acteurs}

\section{Contexte général et motivation du projet}

La \textbf{Faculté} a connu, ces dernières années, une croissance
continue de son parc informatique consécutive à la multiplication
des programmes de formation, à la numérisation de l'enseignement
et à la modernisation des laboratoires de recherche. Cette
croissance s'est accompagnée d'une diversification accrue des
matériels gérés (ordinateurs portables, ordinateurs fixes,
imprimantes laser, imprimantes multifonctions) ainsi que d'une
multiplication des intervenants : départements pédagogiques,
enseignants-chercheurs, fournisseurs accrédités, prestataires de
maintenance externes, etc.

Dans ce contexte, la gestion administrative \emph{ad hoc} qui
prévalait jusqu'ici (tableurs, courriels et registres papier)
montre désormais ses limites : pertes de traçabilité, doublons
d'inventaire, lenteur de traitement des appels d'offres,
difficulté à identifier les fournisseurs récurrents en défaut, et
manque de visibilité sur le taux de panne par catégorie de
matériel. La direction de la Faculté a donc commandité la
conception d'un \textbf{Système d'Information dédié à la gestion
des ressources matérielles}, désigné par le sigle \textbf{S-GRM}
dans la suite du document.

\section{Objectifs stratégiques du projet}

Le projet S-GRM poursuit cinq objectifs stratégiques majeurs,
hiérarchisés ci-dessous par ordre d'importance décroissante~:

\begin{enumerate}[leftmargin=*]
  \item \textbf{Centraliser} l'ensemble des informations relatives
    au parc matériel sur une plateforme unique, accessible aux
    parties prenantes selon des droits différenciés.
  \item \textbf{Tracer} de manière auditable chaque événement
    significatif du cycle de vie d'un équipement : commande,
    livraison, affectation, panne, réparation, mise au rebut.
  \item \textbf{Fluidifier} les processus achats en automatisant
    les étapes répétitives (création d'appels d'offres, comparaison
    chiffrée, blacklistage automatique en cas de récidive).
  \item \textbf{Sécuriser} les transactions et les données
    sensibles (offres financières chiffrées en \mad, identités
    des responsables, historique des décisions).
  \item \textbf{Aider à la décision} en fournissant des
    indicateurs consolidés (budgets consommés, taux de panne,
    note moyenne d'un fournisseur, durée moyenne de réparation).
\end{enumerate}

\section{Périmètre fonctionnel}

Le périmètre couvert par S-GRM se décompose en trois grands
modules transverses, étroitement intégrés mais cohérents sur le
plan métier~:

\begin{itemize}
  \item \textbf{Module Achats (Procurement).} Recueil des besoins
    matériels par département, création d'appels d'offres,
    soumission d'offres par les fournisseurs (date de livraison,
    garantie, prix unitaire et total en \mad), comparaison des
    propositions, sélection du moins-disant qualifié, gestion
    de la liste noire des fournisseurs défaillants.
  \item \textbf{Module Affectation (Allocation).} Catégorisation
    des matériels (ordinateurs, imprimantes), saisie des
    spécifications techniques propres à chaque type, attribution
    d'un code-barres d'inventaire unique, affectation à un
    département ou directement à un enseignant nommément désigné.
  \item \textbf{Module Maintenance.} Signalement des pannes par
    les enseignants utilisateurs, rédaction par les techniciens
    d'un \constat{} comportant explication, date d'apparition,
    fréquence (rare, fréquente, permanente) et type (logiciel ou
    matériel), avec la règle particulière que les pannes
    d'imprimante sont strictement de nature matérielle.
\end{itemize}

Sont \textbf{hors périmètre} de la première version~: la gestion
de la comptabilité générale de la Faculté (qui demeure dans le
système ERP existant), la gestion des stocks de consommables
(toner, papier), et la gestion des contrats de maintenance
externalisés.

\section{Présentation détaillée des acteurs}

Cinq acteurs métier interagissent avec le S-GRM. Chacun dispose
d'un profil de droits spécifique, conçu selon le principe du
\emph{moindre privilège} (\textit{principle of least privilege})
et révisable à tout moment par le Responsable des ressources.

\subsection{Le Responsable des ressources}

Le Responsable des ressources est l'acteur \textbf{pivot} du
système. Il représente la direction administrative de la Faculté
en matière de gestion matérielle.

\paragraph{Droits et responsabilités principales.}
\begin{itemize}
  \item Création, modification et clôture des appels d'offres.
  \item Inscription, modification et suppression des
    fournisseurs dans le référentiel.
  \item Mise sur \textbf{liste noire} d'un fournisseur, sous
    réserve d'apporter un \textbf{motif détaillé} formellement
    enregistré dans le système (par exemple : retards
    systématiques, livraison de matériel non conforme, faillite,
    non-respect de la garantie).
  \item Sélection finale, parmi les offres recevables, du
    \textbf{moins-disant qualifié} dans le respect du règlement
    interne des marchés publics universitaires.
  \item Validation des bons de réception et déclenchement de
    l'affectation des matériels reçus.
  \item Consultation transverse de tous les modules (achats,
    affectation, maintenance) avec accès en lecture intégrale.
\end{itemize}

\paragraph{Limites.} Le Responsable n'a pas vocation à
\emph{rédiger} les constats de panne (rôle du Technicien) ni à
\emph{exprimer} directement les besoins matériels (rôle des Chefs
de département). Cette séparation des responsabilités garantit
l'indépendance hiérarchique et la traçabilité des décisions.

\subsection{Le Chef de département}

Chaque département pédagogique de la Faculté (Informatique,
Mathématiques, Physique, Économie, etc.) est dirigé par un Chef
de département, dont le rôle dans le S-GRM est central pour
l'expression du besoin.

\paragraph{Droits et responsabilités principales.}
\begin{itemize}
  \item Consultation du budget annuel d'investissement matériel
    alloué au département (en \mad).
  \item Saisie et consolidation des besoins matériels remontés
    par les enseignants du département (par exemple~: « 12
    ordinateurs portables i7 16~Go pour la salle TP-3~»).
  \item Validation hiérarchique des demandes avant transmission
    au Responsable des ressources.
  \item Consultation des affectations en cours sur les ressources
    rattachées au département.
  \item Consultation de l'historique des constats de panne
    affectant les ressources du département, à des fins de
    pilotage qualitatif.
\end{itemize}

\paragraph{Limites.} Le Chef de département ne sélectionne pas
les fournisseurs et ne signe pas de bon de commande. Il ne peut
pas consulter les offres d'autres départements. Il n'a pas accès
aux fonctions de blacklistage.

\subsection{L'Enseignant}

L'Enseignant est l'utilisateur final de la grande majorité des
ressources matérielles individuelles (ordinateur de bureau,
ordinateur portable, imprimante personnelle ou partagée).

\paragraph{Droits et responsabilités principales.}
\begin{itemize}
  \item Consultation des ressources qui lui sont nominativement
    affectées (avec leurs spécifications complètes et leur
    code-barres d'inventaire).
  \item \textbf{Signalement d'une panne} sur l'une des ressources
    qui lui sont affectées : description initiale, gravité
    perçue, conditions d'apparition.
  \item Suivi de l'avancement de la prise en charge par les
    techniciens (statut~: signalée, prise en charge, en cours
    de réparation, résolue).
  \item Confirmation de la résolution effective d'une panne après
    intervention, ou ré-ouverture d'un constat en cas de récidive.
\end{itemize}

\paragraph{Limites.} L'Enseignant ne peut pas réaffecter une
ressource à un collègue ; il ne peut ni rédiger ni clôturer un
\constat{} (rôle du Technicien) ; il ne peut pas non plus
exprimer directement un besoin de nouveau matériel auprès du
Responsable~: cette demande remonte par le canal de son Chef de
département.

\subsection{Le Fournisseur}

Le Fournisseur est un acteur \textbf{externe} à la Faculté qui
intervient dans le cadre des procédures d'achat. Il accède au
système par un portail dédié, après inscription validée par le
Responsable des ressources.

\paragraph{Droits et responsabilités principales.}
\begin{itemize}
  \item Inscription en ligne avec saisie obligatoire des champs~:
    raison sociale, ville/localisation, adresse postale complète,
    \textbf{site web} officiel, et \textbf{nom du responsable
    commercial} interlocuteur de la Faculté.
  \item Consultation des appels d'offres ouverts auxquels il est
    autorisé à participer.
  \item \textbf{Soumission d'une offre} pour un appel d'offres,
    comportant obligatoirement~: la \textbf{date de livraison}
    proposée, la \textbf{durée de garantie} en mois, la
    \textbf{marque} et les caractéristiques précises du matériel
    proposé, le \textbf{prix unitaire en \mad}, et le
    \textbf{prix total en \mad} après calcul automatique.
  \item Consultation de l'historique de ses propres soumissions
    et de leur statut (acceptée, rejetée, retenue).
\end{itemize}

\paragraph{Limites.} Un Fournisseur dont le statut est
\textbf{liste noire} se voit refuser toute soumission d'offre
nouvelle de manière automatique par le système. Il peut
néanmoins consulter le motif détaillé de son blacklistage ainsi
que ses anciennes offres archivées. Il ne dispose d'aucun accès
aux modules d'affectation et de maintenance.

\subsection{Le Technicien de maintenance}

Le Technicien de maintenance est l'opérateur technique chargé
d'intervenir physiquement sur les équipements en panne. Il peut
être interne à la Faculté ou rattaché à un prestataire externe
sous contrat.

\paragraph{Droits et responsabilités principales.}
\begin{itemize}
  \item Consultation de la file des pannes signalées,
    triable par urgence, fréquence et type de matériel.
  \item Prise en charge d'un signalement et passage du statut à
    « en cours de diagnostic ».
  \item \textbf{Rédaction du \constat} comportant : explication
    technique détaillée, date d'apparition, fréquence (rare,
    fréquente, permanente), et type (logiciel ou matériel), avec
    la \textbf{règle métier intangible} suivante : pour une
    ressource de type \emph{Imprimante}, la nature de la panne
    est \textbf{exclusivement matérielle}.
  \item Mise à jour de l'état de la ressource (en panne, en
    réparation, hors service, remise en service).
  \item Clôture du \constat{} avec saisie de la solution apportée
    et de la durée d'intervention.
\end{itemize}

\paragraph{Limites.} Le Technicien ne peut ni affecter ni
réaffecter une ressource à un département/enseignant. Il n'a
aucun accès au module Achats. Il ne peut pas modifier
arbitrairement le statut d'un Fournisseur.

\section{Hypothèses, contraintes et exclusions}

\subsection{Hypothèses}
\begin{itemize}
  \item L'authentification des utilisateurs internes (Responsable,
    Chef, Enseignant, Technicien) s'appuie sur l'annuaire
    institutionnel existant (LDAP/Active Directory).
  \item L'authentification des Fournisseurs est gérée localement
    par S-GRM avec mots de passe robustes.
  \item Le réseau interne de la Faculté est considéré comme de
    confiance pour les connexions aux services applicatifs.
\end{itemize}

\subsection{Contraintes techniques}
\begin{itemize}
  \item Toutes les transactions financières sont libellées
    \textbf{exclusivement en Dirhams marocains (\mad)}.
  \item L'interface utilisateur est intégralement \textbf{en
    français}.
  \item L'application doit être déployable sur des serveurs
    \textbf{Linux} via \textbf{Docker}.
\end{itemize}

\subsection{Exclusions}
\begin{itemize}
  \item Pas de gestion comptable analytique.
  \item Pas de gestion des consommables (toner, papier, câbles).
  \item Pas de gestion des bâtiments ni des salles (la salle
    associée à une ressource est un simple champ texte).
\end{itemize}

% ============================================================
% CHAPITRE 2 — EXIGENCES
% ============================================================
\chapter{Exigences fonctionnelles et non-fonctionnelles}

\section{Préambule méthodologique}

Conformément à la norme \emph{IEEE~830-1998} relative à la
spécification d'exigences logicielles, chaque exigence est
identifiée par un code unique de la forme \texttt{EF-Mxx} pour
les exigences fonctionnelles (où \texttt{M} désigne le module et
\texttt{xx} le numéro séquentiel), et \texttt{ENF-yy} pour les
exigences non-fonctionnelles. Chaque exigence est en outre
qualifiée par un niveau de priorité~: \textbf{[B]} bloquante,
\textbf{[E]} essentielle, \textbf{[S]} souhaitable.

\section{Exigences fonctionnelles --- Module Achats (A)}

\subsection{Gestion du référentiel fournisseur}

\begin{description}[leftmargin=*]
  \item[\texttt{EF-A01} {[B]}] Le système doit permettre au
    Responsable d'inscrire un nouveau fournisseur en saisissant,
    au minimum, les champs suivants~: raison sociale, ville,
    adresse postale complète, URL du site web officiel, nom du
    responsable commercial.
  \item[\texttt{EF-A02} {[B]}] Le système doit permettre au
    Responsable de \emph{désactiver} un fournisseur sans le
    supprimer (conservation pour audit).
  \item[\texttt{EF-A03} {[B]}] Le système doit permettre au
    Responsable de placer un fournisseur sur \textbf{liste noire}
    en saisissant un \textbf{motif obligatoire} d'au moins
    cinquante (50) caractères. Toute tentative de blacklistage
    sans motif doit être refusée par le serveur.
  \item[\texttt{EF-A04} {[E]}] Le système doit permettre la
    réhabilitation d'un fournisseur précédemment blacklisté,
    sous réserve de saisir un motif de réhabilitation et la
    signature électronique du Responsable.
\end{description}

\subsection{Cycle de vie d'un appel d'offres}

\begin{description}[leftmargin=*]
  \item[\texttt{EF-A10} {[B]}] Le système doit permettre la
    création d'un appel d'offres avec~: titre, description
    technique, quantité, budget plafond en \mad, date de
    clôture, type de ressource attendue.
  \item[\texttt{EF-A11} {[B]}] Le système doit permettre à un
    fournisseur \emph{actif} de soumettre une offre comportant~:
    date de livraison, durée de garantie en mois, marque, prix
    unitaire en \mad, et prix total en \mad{} (calculé).
  \item[\texttt{EF-A12} {[B]}] Le système doit \textbf{refuser}
    automatiquement toute offre soumise par un fournisseur dont
    le statut courant est \textit{liste noire}, sans recours
    possible.
  \item[\texttt{EF-A13} {[B]}] Le système doit présenter au
    Responsable une vue comparative tabulée des offres
    recevables, triable par prix, garantie ou date de livraison.
  \item[\texttt{EF-A14} {[E]}] Le système doit permettre au
    Responsable de désigner le \textbf{moins-disant qualifié}
    et d'enregistrer dans le journal d'audit le motif de
    sélection.
\end{description}

\section{Exigences fonctionnelles --- Module Affectation (F)}

\begin{description}[leftmargin=*]
  \item[\texttt{EF-F01} {[B]}] Le système doit catégoriser
    chaque ressource selon les types autorisés~: Ordinateur,
    Imprimante.
  \item[\texttt{EF-F02} {[B]}] Pour un Ordinateur, le système
    doit gérer les attributs spécifiques~: marque, processeur
    (CPU), mémoire vive (RAM), disque dur (HDD), écran (taille
    et résolution).
  \item[\texttt{EF-F03} {[B]}] Pour une Imprimante, le système
    doit gérer les attributs spécifiques~: marque, vitesse
    d'impression (pages/minute), résolution (DPI).
  \item[\texttt{EF-F04} {[B]}] Chaque ressource doit recevoir un
    \textbf{numéro d'inventaire / code-barres} unique, généré
    automatiquement par le système et impossible à dupliquer.
  \item[\texttt{EF-F05} {[B]}] Le système doit permettre
    d'affecter une ressource à un \textbf{département} (granularité
    grossière) ou directement à un \textbf{enseignant nommément
    désigné} (granularité fine).
  \item[\texttt{EF-F06} {[E]}] Le système doit conserver
    l'historique complet des affectations passées d'une ressource.
  \item[\texttt{EF-F07} {[S]}] Le système doit permettre l'export
    de l'inventaire au format CSV pour l'audit annuel.
\end{description}

\section{Exigences fonctionnelles --- Module Maintenance (M)}

\begin{description}[leftmargin=*]
  \item[\texttt{EF-M01} {[B]}] Le système doit permettre à un
    Enseignant de signaler une panne sur une ressource qui lui
    est affectée.
  \item[\texttt{EF-M02} {[B]}] Le système doit permettre à un
    Technicien de rédiger un \constat{} avec~: explication,
    date d'apparition, fréquence (\textit{rare},
    \textit{fréquente}, \textit{permanente}), type
    (\textit{logiciel} ou \textit{matériel}).
  \item[\texttt{EF-M03} {[B]}] Le système doit \textbf{rejeter}
    tout \constat{} dont la ressource est de type Imprimante et
    dont le type de panne n'est pas \textit{matériel}.
  \item[\texttt{EF-M04} {[E]}] Le système doit notifier
    automatiquement l'Enseignant initiateur lorsqu'un constat
    le concernant change de statut.
  \item[\texttt{EF-M05} {[E]}] Le système doit permettre de
    clôturer un constat avec saisie de la solution apportée et
    de la durée totale d'intervention.
\end{description}

\section{Exigences non-fonctionnelles}

\subsection{Sécurité (\texttt{ENF-S})}
\begin{description}[leftmargin=*]
  \item[\texttt{ENF-S01}] Toute communication entre le client
    et le serveur doit être chiffrée par TLS~1.2 ou supérieur.
  \item[\texttt{ENF-S02}] Les mots de passe doivent être stockés
    sous forme de condensats salés (algorithme \texttt{bcrypt}
    avec coût $\geq 12$).
  \item[\texttt{ENF-S03}] Le contrôle d'accès doit être basé sur
    les rôles (RBAC), conformément aux profils décrits au
    chapitre 1.
  \item[\texttt{ENF-S04}] Toute action sensible (création
    d'appel d'offre, blacklistage, affectation) doit être
    journalisée dans un journal d'audit infalsifiable.
\end{description}

\subsection{Performance (\texttt{ENF-P})}
\begin{description}[leftmargin=*]
  \item[\texttt{ENF-P01}] Le temps de réponse moyen pour 95~\%
    des requêtes utilisateur ne doit pas dépasser 800~ms.
  \item[\texttt{ENF-P02}] Le système doit pouvoir héberger un
    inventaire de 10\,000 ressources sans dégradation perceptible.
  \item[\texttt{ENF-P03}] Le système doit supporter au moins
    100~utilisateurs concurrents en heure de pointe.
\end{description}

\subsection{Traçabilité (\texttt{ENF-T})}
\begin{description}[leftmargin=*]
  \item[\texttt{ENF-T01}] Chaque enregistrement de la base doit
    porter les colonnes \texttt{date\_creation} et
    \texttt{date\_modification} renseignées automatiquement.
  \item[\texttt{ENF-T02}] Aucun enregistrement métier ne doit
    être supprimé physiquement~: seul le \emph{soft delete}
    (marqueur \texttt{date\_suppression}) est autorisé.
  \item[\texttt{ENF-T03}] Le journal d'audit doit conserver
    l'identité de l'auteur, l'horodatage UTC et la nature de
    l'action effectuée, pour une durée minimale de cinq ans.
\end{description}

\subsection{Disponibilité, ergonomie et maintenabilité}
\begin{description}[leftmargin=*]
  \item[\texttt{ENF-D01}] Disponibilité cible~: 99{,}5~\% sur
    plage 8h--20h en jours ouvrables.
  \item[\texttt{ENF-E01}] L'interface doit respecter les
    standards d'accessibilité WCAG~2.1 niveau AA.
  \item[\texttt{ENF-M01}] Le code source doit présenter un taux
    de couverture de tests automatisés $\geq 70~\%$.
  \item[\texttt{ENF-M02}] Le déploiement doit être entièrement
    reproductible via Docker Compose.
\end{description}

\section{Synthèse des exigences en matrice de criticité}

\begin{longtable}{p{2.5cm} p{8cm} p{1.5cm} p{1.8cm}}
\toprule
\textbf{Code} & \textbf{Intitulé synthétique} & \textbf{Module} &
\textbf{Priorité}\\
\midrule
\endhead
EF-A01 & Inscription d'un fournisseur & Achats & B \\
EF-A03 & Blacklistage avec motif obligatoire & Achats & B \\
EF-A11 & Soumission d'une offre chiffrée & Achats & B \\
EF-A12 & Rejet automatique --- liste noire & Achats & B \\
EF-A14 & Sélection du moins-disant & Achats & E \\
EF-F01 & Catégorisation des ressources & Affectation & B \\
EF-F04 & Code-barres unique d'inventaire & Affectation & B \\
EF-F05 & Affectation à dpt ou enseignant & Affectation & B \\
EF-M01 & Signalement d'une panne & Maintenance & B \\
EF-M02 & Rédaction d'un constat structuré & Maintenance & B \\
EF-M03 & Imprimante $\Rightarrow$ panne matérielle & Maintenance & B \\
ENF-S01 & TLS sur tous les flux & Transverse & B \\
ENF-T03 & Journal d'audit 5 ans & Transverse & E \\
\bottomrule
\caption{Synthèse des exigences principales}
\label{tab:matriceExigences}
\end{longtable}

% ============================================================
% CHAPITRE 3 — PROCESSUS BPMN
% ============================================================
\chapter{Modélisation des processus métier (BPMN)}

\section{Méthodologie}

Les processus métier sont représentés en notation \emph{Business
Process Model and Notation} (BPMN~2.0), produite ici sous forme
de code PlantUML pour permettre une génération automatique des
visuels. Pour chaque processus, nous fournissons une
\emph{introduction longue}, une \emph{description textuelle pas
à pas}, le code PlantUML correspondant, puis un commentaire de
lecture du diagramme.

\section{Processus 1 --- Appel d'offres}

\subsection{Introduction au processus}

Le processus d'\emph{appel d'offres} constitue le cœur du module
Achats. Il s'étend depuis la formalisation administrative du
besoin jusqu'à la notification du fournisseur retenu, en passant
par la phase critique de comparaison et d'élimination des
candidatures non conformes. Trois rôles s'y entrecroisent~: le
Chef de département (qui exprime le besoin), le Responsable des
ressources (qui pilote l'appel d'offres) et le Fournisseur (qui
soumet sa proposition financière).

\subsection{Description textuelle pas à pas}

\begin{enumerate}[leftmargin=*]
  \item Le \textbf{Chef de département} dépose une demande
    matérielle motivée auprès du Responsable.
  \item Le \textbf{Responsable} valide ou refuse la demande.
    En cas de refus, le processus s'arrête.
  \item En cas de validation, le Responsable rédige et publie
    un appel d'offres détaillé.
  \item Les \textbf{Fournisseurs} actifs reçoivent la
    notification et peuvent consulter l'appel.
  \item Chaque Fournisseur intéressé soumet son offre
    (date de livraison, garantie, prix \mad).
  \item Le système \textbf{rejette automatiquement} toute offre
    issue d'un fournisseur en \textit{liste noire}.
  \item À l'issue de la date de clôture, le Responsable
    consulte les offres recevables, les compare et désigne le
    moins-disant qualifié.
  \item Le Responsable notifie le fournisseur retenu et clôt
    l'appel d'offres.
\end{enumerate}

\subsection{Code PlantUML --- Diagramme BPMN}

\begin{lstlisting}[caption={Processus BPMN --- Appel d'offres},
                   label={lst:bpmnAppelOffres}]
@startuml
title Processus BPMN : Cycle de vie d'un appel d'offres

|Chef de departement|
start
:Exprimer un besoin materiel;
:Soumettre la demande au Responsable;

|Responsable des ressources|
:Receptionner la demande;
if (Demande conforme et budget disponible ?) then (oui)
  :Rediger l'appel d'offres;
  :Publier l'appel d'offres;
else (non)
  :Refuser la demande;
  :Notifier le Chef de departement;
  stop
endif

|Fournisseur|
:Consulter les appels d'offres ouverts;
if (Statut fournisseur = liste_noire ?) then (oui)
  :Soumission refusee automatiquement;
  stop
else (non)
  :Saisir une offre (livraison, garantie, prix MAD);
  :Soumettre l'offre;
endif

|Responsable des ressources|
:Recevoir et archiver les offres;
:Comparer les offres recevables;
:Selectionner le moins-disant qualifie;
:Notifier le fournisseur retenu;
:Cloturer l'appel d'offres;

stop
@enduml
\end{lstlisting}

\subsection{Lecture du diagramme}

Le diagramme met en évidence trois \emph{swimlanes} matérialisant
la séparation des responsabilités. Le point de décision sur le
statut du Fournisseur (\texttt{liste\_noire}) apparaît
explicitement comme une \emph{passerelle exclusive}~: c'est ici
que se matérialise l'exigence \texttt{EF-A12}. Toute déviation
sur cette passerelle (par exemple un contournement applicatif)
constituerait une violation directe d'une exigence bloquante.

\section{Processus 2 --- Affectation d'une ressource}

\subsection{Introduction au processus}

L'affectation est le processus par lequel une ressource matérielle
nouvellement acquise (ou démobilisée d'un précédent utilisateur)
est attribuée à une entité utilisatrice. La granularité de
l'affectation peut être grossière (un département entier) ou fine
(un enseignant nommément désigné). Le processus est
systématiquement précédé d'une étape d'inventoriage générant le
code-barres unique imposé par l'exigence \texttt{EF-F04}.

\subsection{Code PlantUML --- Diagramme BPMN}

\begin{lstlisting}[caption={Processus BPMN --- Affectation},
                   label={lst:bpmnAffectation}]
@startuml
title Processus BPMN : Affectation d'une ressource

|Responsable des ressources|
start
:Receptionner physiquement le materiel;
:Saisir les specifications (CPU/RAM/HDD ou vitesse/resolution);
:Generer le code-barres d'inventaire unique;
:Coller l'etiquette code-barres;
:Choisir la cible d'affectation;

if (Affectation a un departement ?) then (oui)
  |Chef de departement|
  :Valider la prise en charge;
  :Repartir l'utilisation au sein du departement;
else (non - enseignant)
  |Enseignant|
  :Signer le bon de prise en charge;
  :Recevoir le materiel;
endif

|Responsable des ressources|
:Enregistrer l'historique d'affectation;
:Cloturer le bon d'affectation;
stop
@enduml
\end{lstlisting}

\subsection{Commentaire}

Le processus se distingue par un \textbf{branchement dual} selon
la cible. La granularité fine (Enseignant) impose une signature
explicite, justifiant la traçabilité forte exigée par
\texttt{ENF-T01}. La génération du code-barres précède
\emph{toujours} l'affectation, garantissant l'unicité et la
traçabilité physique.

\section{Processus 3 --- Maintenance d'une ressource}

\subsection{Introduction}

Le processus de maintenance s'amorce dès qu'un Enseignant
constate un dysfonctionnement. Il se conclut soit par une remise
en service, soit par une mise au rebut (réforme). La règle
métier \texttt{EF-M03} (panne d'imprimante exclusivement
matérielle) trouve son contrôle effectif au moment de la
rédaction du \constat{} par le Technicien.

\subsection{Code PlantUML --- Diagramme BPMN}

\begin{lstlisting}[caption={Processus BPMN --- Maintenance},
                   label={lst:bpmnMaintenance}]
@startuml
title Processus BPMN : Cycle de maintenance

|Enseignant|
start
:Constater un dysfonctionnement;
:Signaler la panne dans le S-GRM;
:Decrire les symptomes observes;

|Technicien de maintenance|
:Recuperer le signalement dans la file;
:Diagnostiquer la ressource;
if (Ressource = Imprimante ?) then (oui)
  :Forcer la nature de panne = MATERIELLE;
else (non)
  :Choisir nature = MATERIELLE ou LOGICIELLE;
endif

:Rediger le Constat de panne;
:Saisir : explication, date, frequence, type;
:Mettre la ressource en etat << en_panne >>;

if (Reparation possible ?) then (oui)
  :Effectuer l'intervention;
  :Tester la remise en service;
  :Cloturer le constat avec solution;
  |Enseignant|
  :Confirmer la remise en service;
else (non)
  |Responsable des ressources|
  :Decider de la reforme;
  :Mettre la ressource en etat << reforme >>;
endif
stop
@enduml
\end{lstlisting}

\subsection{Commentaire}

La passerelle exclusive « Ressource = Imprimante~? » réifie la
règle \texttt{EF-M03}. Un développement défaillant qui omettrait
ce contrôle exposerait la Faculté à des incohérences de données
historiques (statistiques de panne logicielle inappropriées sur
imprimantes), ce qui justifie la mise en place d'un test
d'intégration automatisé spécifique (cf. DAT, chapitre 4).

% ============================================================
% CHAPITRE 4 — CAS D'UTILISATION UML
% ============================================================
\chapter{Modélisation fonctionnelle (UML --- Cas d'utilisation)}

\section{Vue d'ensemble}

Le diagramme de cas d'utilisation présenté ci-après synthétise
l'ensemble des interactions de premier ordre entre les cinq
acteurs et le système. Les relations \emph{include} et
\emph{extend} explicitent les dépendances fonctionnelles
internes. Les cas d'utilisation génériques (authentification,
journalisation) sont mutualisés entre acteurs.

\section{Code PlantUML --- Diagramme de cas d'utilisation global}

\begin{lstlisting}[caption={Diagramme UML --- Cas d'utilisation global},
                   label={lst:umlUseCase}]
@startuml
title Cas d'utilisation : Systeme de Gestion des Ressources Materielles

left to right direction

actor "Responsable\ndes ressources" as Resp
actor "Chef de\ndepartement"      as Chef
actor "Enseignant"                 as Ens
actor "Fournisseur"                as Four
actor "Technicien de\nmaintenance" as Tech

rectangle "Systeme S-GRM" {

  usecase "S'authentifier"               as UC_Auth
  usecase "Consulter le tableau de bord" as UC_Dash

  package "Module Achats" {
    usecase "Inscrire un fournisseur"     as UC_InscFour
    usecase "Mettre un fournisseur\nsur liste noire (avec motif)" as UC_BL
    usecase "Creer un appel d'offres"     as UC_CrAO
    usecase "Soumettre une offre"         as UC_SubOffre
    usecase "Comparer les offres"         as UC_Comp
    usecase "Selectionner le\nmoins-disant" as UC_Selec
  }

  package "Module Affectation" {
    usecase "Cataloguer une ressource"      as UC_Cat
    usecase "Generer un code-barres"        as UC_CB
    usecase "Affecter a un departement"     as UC_AffDept
    usecase "Affecter a un enseignant"      as UC_AffEns
    usecase "Consulter ses ressources"      as UC_Mes
  }

  package "Module Maintenance" {
    usecase "Signaler une panne"               as UC_Sig
    usecase "Rediger un Constat de panne"      as UC_Const
    usecase "Cloturer un Constat"              as UC_CloseC
    usecase "Confirmer la remise en service"   as UC_Conf
  }
}

' Acteurs vers cas d'utilisation
Resp --> UC_InscFour
Resp --> UC_BL
Resp --> UC_CrAO
Resp --> UC_Comp
Resp --> UC_Selec
Resp --> UC_Cat
Resp --> UC_CB
Resp --> UC_AffDept
Resp --> UC_AffEns

Chef --> UC_AffDept
Chef --> UC_Dash

Ens --> UC_Mes
Ens --> UC_Sig
Ens --> UC_Conf

Four --> UC_SubOffre

Tech --> UC_Const
Tech --> UC_CloseC

' Relations
UC_InscFour ..> UC_Auth : <<include>>
UC_BL       ..> UC_Auth : <<include>>
UC_CrAO     ..> UC_Auth : <<include>>
UC_SubOffre ..> UC_Auth : <<include>>
UC_Const    ..> UC_Auth : <<include>>
UC_Sig      ..> UC_Auth : <<include>>

UC_Cat      ..> UC_CB   : <<include>>
UC_Selec    ..> UC_BL   : <<extend>>

@enduml
\end{lstlisting}

\section{Descriptions textuelles longues des cas d'utilisation}

\subsection{UC-01 : Soumettre une offre}
\begin{description}[leftmargin=*]
  \item[Acteur principal] Fournisseur.
  \item[Précondition] Le fournisseur est authentifié et son
    statut courant est \emph{actif} (non blacklisté).
  \item[Scénario nominal]
    \begin{enumerate}
      \item Le fournisseur sélectionne un appel d'offres ouvert.
      \item Il remplit le formulaire d'offre~: date de
        livraison, durée de garantie, marque, prix unitaire en
        \mad, prix total en \mad.
      \item Le système valide la cohérence des champs (prix
        positifs, date de livraison antérieure ou égale à la
        date de clôture acceptée).
      \item Le système enregistre l'offre et notifie le
        Responsable.
    \end{enumerate}
  \item[Scénario alternatif A1] Le fournisseur est blacklisté~:
    le serveur retourne une erreur 400 sans persister l'offre.
  \item[Scénario alternatif A2] Le prix saisi est négatif~: le
    formulaire affiche un message d'erreur de validation.
  \item[Postcondition] L'offre figure dans la liste comparative
    consultable par le Responsable.
\end{description}

\subsection{UC-02 : Mettre un fournisseur sur liste noire}
\begin{description}[leftmargin=*]
  \item[Acteur principal] Responsable des ressources.
  \item[Précondition] Le fournisseur existe dans le référentiel
    et son statut courant est \emph{actif}.
  \item[Scénario nominal]
    \begin{enumerate}
      \item Le Responsable ouvre la fiche du fournisseur
        concerné.
      \item Il sélectionne l'action \emph{Mettre sur liste
        noire}.
      \item Il saisit un motif détaillé (au moins 50 caractères).
      \item Le système enregistre le statut \texttt{liste\_noire}
        et le motif, puis met à jour la date de modification.
      \item Le système informe l'auditeur via le journal
        d'audit.
    \end{enumerate}
  \item[Scénario alternatif] Si le motif est vide ou trop court,
    le système retourne une erreur de validation HTTP 400.
\end{description}

\subsection{UC-03 : Rédiger un constat de panne}
\begin{description}[leftmargin=*]
  \item[Acteur principal] Technicien de maintenance.
  \item[Précondition] Une panne a été signalée par un
    Enseignant et a été prise en charge.
  \item[Scénario nominal]
    \begin{enumerate}
      \item Le Technicien ouvre le signalement.
      \item Il diagnostique la ressource physiquement.
      \item Il sélectionne la nature de panne (logicielle ou
        matérielle).
      \item Il renseigne fréquence, date d'apparition et
        explication.
      \item Le système valide la règle « Imprimante $\Rightarrow$
        matérielle » côté serveur.
      \item Le système persiste le constat et fait passer la
        ressource en état \emph{en\_panne}.
    \end{enumerate}
  \item[Scénario alternatif] Si la ressource est de type
    Imprimante et que la nature saisie est \emph{logicielle}, le
    serveur renvoie HTTP 400 sans persister le constat.
\end{description}

\subsection{UC-04 : Affecter une ressource}
\begin{description}[leftmargin=*]
  \item[Acteur principal] Responsable des ressources.
  \item[Précondition] La ressource est inventoriée
    (code-barres généré) et n'est pas en panne.
  \item[Scénario nominal]
    \begin{enumerate}
      \item Le Responsable sélectionne la ressource.
      \item Il choisit la cible d'affectation~: un département
        ou un enseignant.
      \item Le système crée l'enregistrement d'affectation et
        archive l'éventuelle affectation précédente.
      \item L'enseignant est notifié.
    \end{enumerate}
\end{description}

\subsection{UC-05 : Signaler une panne}
\begin{description}[leftmargin=*]
  \item[Acteur principal] Enseignant.
  \item[Précondition] L'enseignant est affecté à la ressource.
  \item[Scénario nominal]
    \begin{enumerate}
      \item L'enseignant sélectionne sa ressource en panne.
      \item Il décrit le symptôme constaté.
      \item Il indique une éventuelle gravité ressentie.
      \item Le système enregistre le signalement et notifie
        les techniciens.
    \end{enumerate}
\end{description}

\part{Document d'Architecture Technique}

% ============================================================
% CHAPITRE 1 — Présentation de l'architecture
% ============================================================
\chapter{Présentation et justification de l'architecture}

\section{Vue d'ensemble : architecture trois-tiers}

L'architecture retenue pour S-GRM est une \textbf{architecture
trois-tiers} (\emph{three-tier architecture}), modèle de
référence pour les applications web modernes. Cette segmentation
horizontale en trois couches répond directement à plusieurs des
exigences non-fonctionnelles formulées au chapitre 2 du cahier
des charges, en particulier la maintenabilité (\texttt{ENF-M01})
et la performance (\texttt{ENF-P02}).

\begin{itemize}
  \item \textbf{Tier 1 --- Présentation.} Interface utilisateur
    en \textbf{React}, exécutée dans le navigateur. Elle prend
    en charge le rendu, l'ergonomie, la validation
    intermédiaire des saisies et l'expérience utilisateur (UX).
  \item \textbf{Tier 2 --- Logique applicative.} Serveur
    \textbf{Node.js / Express.js} exposant une API REST. Il
    embarque la totalité des règles de gestion (validation des
    blacklists, contrôle imprimante = matérielle, calculs de
    prix, etc.) et orchestre les transactions vers la base de
    données.
  \item \textbf{Tier 3 --- Persistance.} Base de données
    relationnelle \textbf{PostgreSQL}, choisie pour la richesse
    de son SQL et la robustesse de ses garanties ACID.
\end{itemize}

\section{Justification du choix de la stack PERN}

Le sigle \textbf{PERN} désigne l'association
\textbf{P}ostgreSQL, \textbf{E}xpress, \textbf{R}eact,
\textbf{N}ode.js. Plusieurs raisons motivent ce choix dans le
contexte académique de la Faculté.

\subsection{Pourquoi PostgreSQL ?}

PostgreSQL est aujourd'hui considéré comme le système de gestion
de base de données relationnelles open source le plus avancé.
Son adoption pour le projet repose sur cinq piliers~:

\begin{enumerate}
  \item \textbf{Transactionnel ACID complet.} L'écriture d'un
    \constat{} et la mise à jour conjointe de l'état de la
    ressource doivent être atomiques pour éviter les
    incohérences. PostgreSQL offre un support transactionnel de
    référence (sérialisable, lecture consistante des snapshots).
  \item \textbf{Types riches.} Les colonnes \texttt{JSONB} sont
    utilisées pour stocker les spécifications techniques
    polymorphes des ressources (CPU/RAM/HDD pour un ordinateur,
    vitesse/résolution pour une imprimante) sans recourir à une
    table fille par sous-classe, tout en conservant
    l'indexabilité.
  \item \textbf{Contraintes intégrées.} Les types \texttt{ENUM}
    permettent de matérialiser au niveau de la base les statuts
    \texttt{actif/liste\_noire} et \texttt{disponible/en\_panne/
    \ldots}, défendant ainsi les invariants même en cas de bug
    applicatif.
  \item \textbf{Open source et libre.} Aucune redevance, ce qui
    convient parfaitement à un budget académique.
  \item \textbf{Écosystème mature.} Excellents pilotes Node
    (\texttt{pg}, \texttt{sequelize}), nombreux outils
    d'administration (\texttt{pgAdmin}), procédures de
    sauvegarde éprouvées.
\end{enumerate}

\subsection{Pourquoi Express.js ?}

Express est le \emph{de facto} standard pour bâtir des API REST
en Node.js. Sa philosophie minimaliste laisse libre le choix des
middlewares (CORS, parsing JSON, journalisation), ce qui
favorise une compréhension fine de la pile par les étudiants.
Son modèle d'intergiciels enchaînés (\emph{middleware pipeline})
se prête particulièrement bien à l'application transverse de
règles non-fonctionnelles : authentification, journal d'audit,
limitation de débit.

\subsection{Pourquoi React ?}

React, avec ses \emph{hooks} (\texttt{useState},
\texttt{useEffect}, \texttt{useCallback}), permet une
modélisation déclarative et composable de l'interface. Le pattern
de composants fonctionnels purs facilite la réutilisation et le
test unitaire (par exemple via \texttt{Jest} et \texttt{React
Testing Library}). Pour le S-GRM, la décomposition naturelle est
\texttt{InventaireRessources}, \texttt{ListeFournisseurs},
\texttt{FormulaireRessource}, etc., avec un service \texttt{api}
mutualisé.

\subsection{Pourquoi Node.js ?}

Le partage du langage JavaScript entre frontend et backend
réduit la charge cognitive de l'équipe, mutualise certains
utilitaires (validation, formats de date) et facilite la
montée en compétences. Le modèle d'événements non bloquants de
Node convient à un système majoritairement IO-bound (lectures
BD, appels HTTP).

\subsection{Cohérence transverse de la stack}

L'unification du langage, la portabilité Linux/Windows/macOS,
l'écosystème npm très riche et la conteneurisation aisée via
Docker rendent PERN homogène, raisonnablement performante,
maintenable, et bien adaptée à un projet académique de cette
envergure.

\section{Patrons d'architecture employés}

\begin{description}[leftmargin=*]
  \item[Modèle--Vue--Contrôleur (MVC).] La couche Express
    applique strictement la séparation Routes $\rightarrow$
    Contrôleurs $\rightarrow$ Services $\rightarrow$ Dépôts
    $\rightarrow$ Modèles, qui évite tout couplage entre
    transport HTTP et logique métier.
  \item[Repository.] L'accès aux données est encapsulé dans des
    classes \texttt{DepotXxx}, ce qui permet de changer
    ultérieurement la couche de persistance sans toucher aux
    services.
  \item[Factory Method.] La fabrique \texttt{FabriqueRessource}
    construit les bons objets selon le type (Ordinateur ou
    Imprimante), encapsulant les validations propres à chaque
    sous-type.
  \item[Singleton.] La classe \texttt{BaseDeDonnees} expose une
    instance unique de \texttt{Sequelize} pour éviter la
    multiplication des connexions.
\end{description}

% ============================================================
% CHAPITRE 2 — Modélisation des données
% ============================================================
\chapter{Modélisation des données}

\section{Diagramme de classes UML}

\begin{lstlisting}[caption={Diagramme UML --- Classes du domaine},
                   label={lst:umlClasses}]
@startuml
title Diagramme de classes : domaine S-GRM

abstract class Ressource {
  +id : int
  +nom : string
  +marque : string
  +prix_unitaire_mad : decimal
  +etat : EtatRessource
  +date_creation : datetime
  +code_barres : string
  --
  +afficherFiche() : void
  +affecter(cible) : void
  +mettreEnPanne() : void
}

class Ordinateur {
  +cpu : string
  +ram : string
  +hdd : string
  +ecran : string
}

class Imprimante {
  +vitesse_impression : string
  +resolution : string
}

class Fournisseur {
  +id : int
  +nom : string
  +ville : string
  +adresse : string
  +site_web : string
  +nom_responsable : string
  +telephone : string
  +statut : StatutFournisseur
  +motif_liste_noire : string
  --
  +peutSoumettreOffre() : bool
  +mettreEnListeNoire(motif) : void
}

class AppelOffres {
  +id : int
  +titre : string
  +description : string
  +budget_max_mad : decimal
  +date_cloture : date
  +statut : StatutAO
  --
  +cloturer() : void
  +selectionnerMoinsDisant() : Offre
}

class Offre {
  +id : int
  +date_livraison : date
  +garantie_mois : int
  +marque : string
  +prix_unitaire_mad : decimal
  +prix_total_mad : decimal
  +statut : StatutOffre
}

class ConstatDePanne {
  +id : int
  +date_constat : date
  +nature_panne : NaturePanne
  +description : text
  +frequence : Frequence
  +urgence : Urgence
  +solution : text
}

enum StatutFournisseur {
  ACTIF
  LISTE_NOIRE
}

enum EtatRessource {
  DISPONIBLE
  EN_SERVICE
  EN_PANNE
  REFORME
}

enum NaturePanne {
  MATERIELLE
  LOGICIELLE
}

enum Frequence {
  RARE
  FREQUENTE
  PERMANENTE
}

' Heritage
Ressource <|-- Ordinateur
Ressource <|-- Imprimante

' Associations
Fournisseur "1" -- "0..*" Ressource : fournit >
Fournisseur "1" -- "0..*" Offre     : soumet >
AppelOffres "1" -- "0..*" Offre     : recoit >
Ressource   "1" -- "0..*" ConstatDePanne : subit >
@enduml
\end{lstlisting}

\section{Dictionnaire de données}

\subsection{Entité \texttt{Ressource} (abstraite)}

\begin{longtable}{p{4cm} p{3.5cm} p{6cm}}
\toprule
\textbf{Attribut} & \textbf{Type} & \textbf{Description} \\
\midrule
\endhead
\texttt{id} & SERIAL & Clé primaire auto-incrémentée. \\
\texttt{nom} & VARCHAR(255) & Libellé de la ressource. \\
\texttt{marque} & VARCHAR(255) & Fabricant. \\
\texttt{prix\_unitaire\_mad} & DECIMAL(12,2) & Prix unitaire en
\mad{} (toujours positif). \\
\texttt{etat} & ENUM & disponible / en\_service / en\_panne /
reforme. \\
\texttt{specifications} & JSONB & Spécifications techniques
polymorphes selon le type. \\
\texttt{fournisseur\_id} & INT & FK vers \texttt{fournisseurs.id}. \\
\texttt{code\_barres} & VARCHAR(64) & Code-barres unique. \\
\texttt{date\_creation} & TIMESTAMP & Date d'inventoriage. \\
\texttt{date\_modification} & TIMESTAMP & Date de dernière mise à jour. \\
\bottomrule
\caption{Dictionnaire --- Ressource}
\label{tab:dictRessource}
\end{longtable}

\subsection{Entité \texttt{Fournisseur}}

\begin{longtable}{p{4cm} p{3.5cm} p{6cm}}
\toprule
\textbf{Attribut} & \textbf{Type} & \textbf{Description} \\
\midrule
\endhead
\texttt{id} & SERIAL & Clé primaire. \\
\texttt{nom} & VARCHAR(255) & Raison sociale. \\
\texttt{ville} & VARCHAR(120) & Localisation principale. \\
\texttt{adresse} & TEXT & Adresse postale complète. \\
\texttt{site\_web} & VARCHAR(255) & URL officielle. \\
\texttt{nom\_responsable} & VARCHAR(255) & Interlocuteur. \\
\texttt{telephone} & VARCHAR(20) & Numéro de téléphone. \\
\texttt{statut} & ENUM & actif / liste\_noire. \\
\texttt{motif\_liste\_noire} & TEXT & Justification obligatoire
si statut = liste\_noire. \\
\texttt{date\_creation} & TIMESTAMP & Inscription. \\
\texttt{date\_modification} & TIMESTAMP & Mise à jour. \\
\bottomrule
\caption{Dictionnaire --- Fournisseur}
\label{tab:dictFournisseur}
\end{longtable}

\subsection{Entité \texttt{ConstatDePanne}}

\begin{longtable}{p{4cm} p{3.5cm} p{6cm}}
\toprule
\textbf{Attribut} & \textbf{Type} & \textbf{Description} \\
\midrule
\endhead
\texttt{id} & SERIAL & Clé primaire. \\
\texttt{ressource\_id} & INT & FK vers \texttt{ressources.id}. \\
\texttt{date\_constat} & DATE & Date d'apparition. \\
\texttt{nature\_panne} & ENUM & matérielle / logicielle (forcée
matérielle si Imprimante). \\
\texttt{description} & TEXT & Explication détaillée. \\
\texttt{frequence} & ENUM & rare / fréquente / permanente. \\
\texttt{urgence} & ENUM & basse / moyenne / haute / critique. \\
\texttt{solution} & TEXT & Action correctrice apportée. \\
\texttt{date\_creation} & TIMESTAMP & Date de création de la
trace. \\
\bottomrule
\caption{Dictionnaire --- Constat de Panne}
\label{tab:dictConstat}
\end{longtable}

\subsection{Relations principales}

\begin{itemize}
  \item Un \fournisseur{} \textbf{fournit} 0..* \ressource{}.
  \item Un \fournisseur{} \textbf{soumet} 0..* \emph{Offres}.
  \item Une \ressource{} \textbf{subit} 0..* \constat{}.
  \item Un \emph{Appel d'offres} \textbf{reçoit} 0..* \emph{Offres}.
  \item L'héritage \textbf{Ressource $\leftarrow$ Ordinateur},
    \textbf{Ressource $\leftarrow$ Imprimante} est implémenté
    par discriminateur \texttt{type} + \texttt{specifications JSONB}
    (héritage de table simple).
\end{itemize}

% ============================================================
% CHAPITRE 3 — Architectures physique et logique
% ============================================================
\chapter{Architectures physique et logique}

\section{Diagramme de composants}

\begin{lstlisting}[caption={Diagramme UML --- Composants},
                   label={lst:umlComposants}]
@startuml
title Diagramme de composants : interaction des couches

package "Navigateur (Client)" {
  [Composants React] as React
  [Service apiService.js] as ApiSvc
}

package "Serveur applicatif" {
  [Express.js] as Express
  [Controleurs] as Ctrl
  [Services metier] as Svc
  [Depots Sequelize] as Repo
}

database "PostgreSQL" as PG

React --> ApiSvc : appels fetch/JSON
ApiSvc --> Express : HTTPS/REST (JSON)
Express --> Ctrl
Ctrl --> Svc : delegation logique
Svc --> Repo : appels CRUD
Repo --> PG : SQL via Sequelize ORM

@enduml
\end{lstlisting}

\subsection{Commentaire}

Le composant \texttt{Service apiService.js} matérialise la
\emph{façade} HTTP côté client~: il centralise les appels
\texttt{fetch}, normalise le format des réponses
(\texttt{\{succes, donnees, message\}}) et abstrait les composants
React de la complexité du protocole.

\section{Diagramme de déploiement}

\begin{lstlisting}[caption={Diagramme UML --- Déploiement Docker},
                   label={lst:umlDeploiement}]
@startuml
title Diagramme de deploiement : conteneurs Docker

node "Hote Docker (Linux)" as Host {

  node "Conteneur grm_frontend" {
    component "Nginx (servant React build)" as Nginx
  }

  node "Conteneur grm_backend" {
    component "Node.js + Express" as Node
  }

  node "Conteneur grm_postgres" {
    database "PostgreSQL 16" as PG
    folder "Volume grm_donnees_postgres" as Vol
  }

  Nginx --> Node : proxy /api -> backend:3001
  Node --> PG : SQL (Sequelize) sur reseau bridge
  PG --> Vol : persistance des donnees
}

actor "Utilisateur (navigateur)" as User
User --> Nginx : HTTP/80
@enduml
\end{lstlisting}

\subsection{Commentaire}

Trois conteneurs orchestrés par \texttt{docker-compose}~: le
frontend Nginx (port 80, qui sert les fichiers statiques React
et joue le rôle de \emph{reverse proxy} vers le backend), le
backend Express (port interne 3001) et la base PostgreSQL
(volume nommé \texttt{grm\_donnees\_postgres} pour la
persistance entre redémarrages). Tous les conteneurs partagent
un réseau Docker bridge \texttt{reseau\_grm}, ce qui leur
permet de se résoudre mutuellement par leur nom de service.

\section{Diagramme de packages (organisation du code source)}

\begin{lstlisting}[caption={Diagramme UML --- Packages backend},
                   label={lst:umlPackages}]
@startuml
title Organisation du code source backend

package "src" {
  package "config" {
    [BaseDeDonnees.js]
  }
  package "models" {
    [Ressource.js]
    [Fournisseur.js]
    [ConstatDePanne.js]
    [index.js]
  }
  package "repositories" {
    [DepotDeBase.js]
    [DepotRessource.js]
    [DepotFournisseur.js]
    [DepotConstatDePanne.js]
  }
  package "factories" {
    [FabriqueRessource.js]
  }
  package "services" {
    [ServiceRessource.js]
    [ServiceFournisseur.js]
    [ServiceConstatDePanne.js]
  }
  package "controllers" {
    [ControleurRessource.js]
    [ControleurFournisseur.js]
    [ControleurConstatDePanne.js]
  }
  package "routes" {
    [routesRessource.js]
    [routesFournisseur.js]
    [routesConstatDePanne.js]
  }
  [app.js]
}

[app.js] --> [routes]
[routes] --> [controllers]
[controllers] --> [services]
[services] --> [repositories]
[services] --> [factories]
[repositories] --> [models]
[models] --> [config]
@enduml
\end{lstlisting}

% ============================================================
% CHAPITRE 4 — Comportement système
% ============================================================
\chapter{Comportement système : diagramme de séquence}

\section{Scénario : un fournisseur soumet une offre chiffrée en MAD}

\subsection{Code PlantUML --- Diagramme de séquence}

\begin{lstlisting}[caption={Séquence --- Soumission d'une offre par un fournisseur},
                   label={lst:umlSequence}]
@startuml
title Sequence : soumission d'une offre chiffree en MAD

actor Fournisseur as F
boundary "React\n(FormulaireOffre)" as UI
control "apiService.js" as Api
participant "Express\n(Routeur)"     as Rt
participant "ControleurFournisseur"  as Ctrl
participant "ServiceFournisseur"     as Svc
participant "DepotFournisseur"       as Dep
database   "PostgreSQL"              as DB

F -> UI : Saisir prix, livraison, garantie
UI -> Api : creerOffre(payload JSON, MAD)
Api -> Rt : POST /api/offres (JSON, TLS)

Rt -> Ctrl : invoquerControleur(req, res)
Ctrl -> Svc : validerEtCreerOffre(payload)

Svc -> Dep : trouverParId(fournisseurId)
Dep -> DB : SELECT * FROM fournisseurs WHERE id = ?
DB --> Dep : ligne fournisseur
Dep --> Svc : objet Fournisseur

alt Statut = liste_noire
  Svc --> Ctrl : throw "Offre rejetee : liste noire"
  Ctrl --> Rt : 400 Bad Request
  Rt --> Api : reponse JSON {succes:false, message}
  Api --> UI : afficher message d'erreur
  UI --> F : "Votre societe est sur liste noire."
else Statut = actif
  Svc -> Dep : creer(offre)
  Dep -> DB : BEGIN; INSERT INTO offres (...) VALUES (...);
  DB --> Dep : id genere
  Dep --> Svc : Offre persistee
  Svc -> Dep : COMMIT
  Svc --> Ctrl : retour offre
  Ctrl --> Rt : 201 Created + JSON
  Rt --> Api : reponse {succes:true, donnees}
  Api --> UI : confirmer creation
  UI --> F : "Offre soumise avec succes."
end
@enduml
\end{lstlisting}

\subsection{Explication technique pas à pas}

\begin{enumerate}[leftmargin=*]
  \item \textbf{Saisie côté client.} Le composant React
    \texttt{FormulaireOffre} maintient un état local via
    \texttt{useState}. À la soumission, il appelle la fonction
    \texttt{creerOffre} du service \texttt{apiService.js}.
  \item \textbf{Préparation de la requête HTTP.} Le service
    construit un objet JSON contenant les champs de l'offre, en
    s'assurant que \texttt{prix\_unitaire\_mad} est un nombre,
    puis émet une requête \texttt{POST /api/offres} avec en-tête
    \texttt{Content-Type: application/json}. Le canal est
    sécurisé par TLS conformément à \texttt{ENF-S01}.
  \item \textbf{Routage Express.} Express consomme le corps via
    le middleware \texttt{express.json()}, applique CORS, puis
    aiguille la requête vers le routeur \texttt{routesFournisseur}
    qui invoque la fonction de contrôleur appropriée.
  \item \textbf{Délégation au service.} Le contrôleur effectue
    une validation minimale (présence des champs obligatoires)
    et délègue la logique métier au \texttt{ServiceFournisseur}.
  \item \textbf{Vérification du statut.} Le service interroge le
    \texttt{DepotFournisseur} qui exécute une requête
    \texttt{SELECT} ciblée sur la clé primaire. Si le statut est
    \texttt{liste\_noire}, une exception est levée et propagée
    en HTTP 400 (sans persistance de l'offre).
  \item \textbf{Insertion transactionnelle.} Si le statut est
    \texttt{actif}, le service ouvre une transaction (BEGIN),
    insère l'offre, puis valide (COMMIT). Sequelize gère ces
    appels via une connexion poolée.
  \item \textbf{Réponse au client.} Le contrôleur renvoie
    \texttt{201 Created} avec un payload JSON normalisé. La
    couche React met à jour l'interface et affiche un toast de
    confirmation à l'utilisateur.
  \item \textbf{Journalisation d'audit.} Un middleware annexe
    enregistre l'événement (qui, quand, quoi) dans la table
    \texttt{journal\_audit}, conformément à l'exigence
    \texttt{ENF-T03}.
\end{enumerate}

\section{Conclusion}

Le présent document a couvert l'intégralité du périmètre
fonctionnel et technique du Système de Gestion des Ressources
Matérielles. Les choix d'architecture — trois-tiers, stack PERN,
patrons MVC/Repository/Factory/Singleton — répondent
rigoureusement aux exigences fonctionnelles et non-fonctionnelles
formulées par la maîtrise d'ouvrage. L'ensemble des diagrammes
PlantUML fournis sont directement intégrables à un environnement
de génération automatique de documentation (par exemple via
\texttt{plantuml.jar}). Les futures itérations pourront porter
sur l'authentification fédérée SSO, la dématérialisation
complète des bons de commande et l'intégration d'un module
décisionnel (BI) couplé aux indicateurs de pannes.

\appendix
\chapter*{Annexe A --- Glossaire}
\addcontentsline{toc}{chapter}{Annexe A --- Glossaire}

\begin{description}[leftmargin=*]
  \item[ACID] Atomicité, Cohérence, Isolation, Durabilité.
  \item[BPMN] Business Process Model and Notation.
  \item[CRUD] Create, Read, Update, Delete.
  \item[MAD] Dirham marocain, monnaie officielle.
  \item[MVC] Modèle--Vue--Contrôleur.
  \item[ORM] Object-Relational Mapping.
  \item[PERN] PostgreSQL + Express + React + Node.
  \item[REST] Representational State Transfer.
  \item[S-GRM] Système de Gestion des Ressources Matérielles.
  \item[SSO] Single Sign-On.
  \item[TLS] Transport Layer Security.
\end{description}

\end{document}
